\DocumentMetadata{
  pdfversion=2.0,
  pdfstandard=ua-2,
  lang=en-US,
  tagging=on,
  tagging-setup={math/setup=mathml-SE,tabsorder=structure}
}
\documentclass[11pt,letterpaper]{article}
\usepackage[margin=0.68in,includefoot]{geometry}
\usepackage{newcomputermodern}
\usepackage{amsmath,amscd,mathtools}
\usepackage{tikz,adjustbox}
\providecommand{\square}{\mdlgwhtsquare}
\usepackage[hidelinks]{hyperref}
\hypersetup{
  pdftitle={Analysis PhD Exam, May 2000},
  pdfauthor={University of Florida Department of Mathematics},
  pdfsubject={Graduate mathematics examination},
  pdfdisplaydoctitle=true,
  bookmarksopen=true
}
\setcounter{secnumdepth}{0}
\setlength{\parindent}{0pt}
\setlength{\parskip}{0.28em}
\setlength{\emergencystretch}{3em}
\setlength{\leftmargini}{2.15em}
\setlength{\leftmarginii}{2.65em}
\setlength{\leftmarginiii}{3em}
\renewcommand{\labelenumi}{\textbf{\theenumi.}}
\pagestyle{empty}
\raggedbottom
\allowdisplaybreaks
\newenvironment{examproblems}[1]{%
  \begin{enumerate}%
  \setcounter{enumi}{#1}%
  \setlength{\topsep}{0.35em}%
  \setlength{\partopsep}{0pt}%
  \setlength{\itemsep}{0.48em}%
  \setlength{\parsep}{0.18em}%
}{\end{enumerate}}
\newenvironment{examsubparts}{%
  \begin{enumerate}%
  \setlength{\topsep}{0.18em}%
  \setlength{\partopsep}{0pt}%
  \setlength{\itemsep}{0.16em}%
  \setlength{\parsep}{0.1em}%
}{\end{enumerate}}
\newenvironment{examinstructions}{%
  \par\begingroup\itshape
}{%
  \par\endgroup\vspace{0.18em}
}
\makeatletter
\renewcommand\section{\@startsection{section}{1}{0pt}%
  {-1.25ex plus -.25ex minus -.1ex}%
  {0.55ex plus .1ex}%
  {\normalfont\large\bfseries}}
\renewcommand{\@maketitle}{%
  \newpage\null\vskip -1em%
  {\footnotesize\bfseries
    \href{https://www.ufl.edu/}{University of Florida}, \href{https://math.ufl.edu/}{Department of Mathematics}\par}%
  \vskip -0.5em%
  \tagstructbegin{tag=Title}%
    {\LARGE\bfseries\href{https://gma.math.ufl.edu/exams/analysis-phd/}{Analysis PhD Exam}, May 2000\par}%
  \tagstructend%
  \vskip 0.25em%
  \tagmcbegin{artifact}%
    \hrule height 1pt%
  \tagmcend%
  \vskip 1em%
}
\makeatother
\title{Analysis PhD Exam, May 2000}
\author{}
\date{}
\begin{document}
\maketitle
\thispagestyle{empty}
\begin{examinstructions}
Be sure to write all proofs in a neat and logical fashion. Give reasons for all steps!
\end{examinstructions}
\begin{examproblems}{0}
\item
State and prove the Vitali Convergence Theorem for integrals.
\item
State and prove the Martingale Convergence Theorem.
\item
Let \((X,\mathcal B,\mu)\) and \((Y,\mathcal J,\nu)\) be finite measure spaces. Construct \(\mu\otimes\nu\) on \(\mathcal B\otimes\mathcal J\). Give details.
\item
Let \(\{\mu_m\}\) be a sequence of real valued signed measures on \((S,\Sigma)\). Define 
\[
\lambda(A)=\sum_{m=1}^{\infty}\frac{1}{2^m}\frac{|\mu_m|(A)}{1+|\mu_m|(S)},\qquad A\in\Sigma.
\]
 Show~\(\lambda\) is a finite measure on~\(\Sigma\) and \(\mu_m\overset{\varepsilon-\delta}{\ll}\lambda\) for each~\(m\).
\item
Let \(f:[0,b]\to\mathbb R\) be Lebesgue integrable. Suppose \(\int_{[0,c]}f\,dm=0\) for each \(c\in[0,b]\). Show \(f=0\) a.e. \(m\).
\item
Let \((\Omega,\mathcal F,P)\) be a pr. space, \(Y\) an integrable r.v. and \(\mathcal E\) a sub σ-field of \(\mathcal F\). Suppose \(\mathcal E\) and \(\mathcal F(Y)\) are independent. What is \(E(Y\mid\mathcal E)\)? Prove.
\item
Let \(\mathcal X\) be a B-space and suppose \(\mathcal Y\) and \(\mathcal Z\) are closed linear subspaces of \(\mathcal X\) such that each \(x\in\mathcal X\) has a unique representation of the form \(x=y+z\), where \(y\in\mathcal Y\) and \(z\in\mathcal Z\). Show that there is a constant~\(K\) such that \(|y|\le K|x|\) and \(|z|\le K|x|\), for each \(x\in\mathcal X\), where~\(y\) \& \(z\) are given as above.

\par
Hint: define a map \(T:\mathcal X\to\mathcal Y\) which is useful, use the closed graph theorem to get continuity; then define \(S:\mathcal X\to\mathcal Z\), etc.
\end{examproblems}
\end{document}
